\documentclass[12pt,a4paper]{article}

% ------------------------------------------------
% PAQUETES
% ------------------------------------------------

\usepackage[spanish, es-nodecimaldot]{babel}
\usepackage[utf8]{inputenc}
\usepackage[T1]{fontenc}

\usepackage{amsmath}
\usepackage{amssymb}
\usepackage{amsfonts}
\usepackage{mathtools}

\usepackage{geometry}
\geometry{
    margin=2.5cm
}

\usepackage{titlesec}
\usepackage{xcolor}

\usepackage{tcolorbox}
\tcbuselibrary{breakable}

\usepackage{enumitem}

\usepackage{setspace}
\setstretch{1.15}

% ------------------------------------------------
% ESTILO
% ------------------------------------------------

\definecolor{mainblue}{RGB}{25,55,109}
\definecolor{lightblue}{RGB}{232,240,255}

\titleformat{\section}
{\normalfont\Large\bfseries\color{mainblue}}
{\thesection.}{0.5em}{}

\titleformat{\subsection}
{\normalfont\large\bfseries\color{mainblue}}
{}{0em}{}

% Caja para identidades
\newtcolorbox{identidad}{
    breakable,
    colback=lightblue,
    colframe=mainblue,
    boxrule=0.8pt,
    arc=3mm,
    left=3mm,
    right=3mm,
    top=2mm,
    bottom=2mm
}

% Caja para pasos
\newtcolorbox{paso}{
    breakable,
    colback=white,
    colframe=black!40,
    boxrule=0.5pt,
    arc=2mm,
    left=2mm,
    right=2mm,
    top=1mm,
    bottom=1mm
}

% ------------------------------------------------
% DOCUMENTO
% ------------------------------------------------

\begin{document}

\begin{center}

{\Huge \textbf{Solucionario}} \\[0.3cm]

{\LARGE Identidades trigonométricas}

\vspace{0.5cm}

\rule{0.9\textwidth}{0.5pt}

\end{center}

\vspace{0.8cm}

\noindent
\textbf{Nota:} En cada ejercicio denotamos con \(L\) la expresión del lado izquierdo de la igualdad. 
Iremos transformando \(L\) paso a paso, utilizando únicamente identidades trigonométricas y álgebra, 
hasta obtener exactamente la expresión del otro lado de la igualdad. Nunca usaremos el otro lado de la igualdad durante la simplificación; solo al final verificamos que la expresión obtenida coincide con él. 
En algunos casos hacemos un pequeño cambio de variable (por ejemplo, \(t = \tan\alpha\)) 
para facilitar la manipulación, siempre manteniendo la equivalencia.



% =====================================================
% a)
% =====================================================
\section*{a)}

\begin{identidad}
\[
\frac{\csc\alpha}{1+\csc\alpha}
-
\frac1{\sin\alpha-1}
=
2\sec^2\alpha
\]
\end{identidad}

\subsection*{Paso 1}
Partimos del lado izquierdo:
\[
L = \frac{\csc\alpha}{1+\csc\alpha} - \frac1{\sin\alpha-1}
\]
Escribimos \(\csc\alpha = \frac1{\sin\alpha}\):
\[
L = \frac{\frac1{\sin\alpha}}{1+\frac1{\sin\alpha}} - \frac1{\sin\alpha-1}
\]
Multiplicamos numerador y denominador de la primera fracción por \(\sin\alpha\):
\[
L = \frac1{1+\sin\alpha} - \frac1{\sin\alpha-1}
\]

\subsection*{Paso 2}
Observamos que \(\sin\alpha-1 = -(1-\sin\alpha)\), luego:
\[
-\frac1{\sin\alpha-1} = \frac1{1-\sin\alpha}
\]
Por tanto:
\[
L = \frac1{1+\sin\alpha} + \frac1{1-\sin\alpha}
\]

\subsection*{Paso 3}
Sumamos las fracciones con denominador común \((1+\sin\alpha)(1-\sin\alpha)\):
\[
L = \frac{(1-\sin\alpha)+(1+\sin\alpha)}{(1+\sin\alpha)(1-\sin\alpha)}
= \frac{2}{1-\sin^2\alpha}
\]
Usamos \(1-\sin^2\alpha = \cos^2\alpha\):
\[
L = \frac{2}{\cos^2\alpha} = 2\sec^2\alpha
\]
Hemos obtenido el lado derecho.

\newpage
% =====================================================
% b)
% =====================================================
\section*{b)}

\begin{identidad}
\[
\frac{\tan\alpha-\cot\alpha}
{1-2\cos^2\alpha}
=
\tan\alpha+\cot\alpha
\]
\end{identidad}

\subsection*{Paso 1}
Partimos del lado izquierdo:
\[
L = \frac{\tan\alpha-\cot\alpha}{1-2\cos^2\alpha}
\]
Escribimos en seno y coseno:
\[
\tan\alpha-\cot\alpha = \frac{\sin\alpha}{\cos\alpha} - \frac{\cos\alpha}{\sin\alpha}
= \frac{\sin^2\alpha-\cos^2\alpha}{\sin\alpha\cos\alpha}
\]
Además, \(1-2\cos^2\alpha = - (2\cos^2\alpha-1)\). Notamos que \(\sin^2\alpha-\cos^2\alpha = -(\cos^2\alpha-\sin^2\alpha) = -(2\cos^2\alpha-1)\).

\subsection*{Paso 2}
Sustituimos:
\[
L = \frac{ \frac{\sin^2\alpha-\cos^2\alpha}{\sin\alpha\cos\alpha} }{1-2\cos^2\alpha}
= \frac{ \frac{-(2\cos^2\alpha-1)}{\sin\alpha\cos\alpha} }{ -(2\cos^2\alpha-1) }
\]
Cancelamos \(-(2\cos^2\alpha-1)\) (suponiendo que no es cero):
\[
L = \frac{1}{\sin\alpha\cos\alpha}
\]

\subsection*{Paso 3}
Transformamos \(\frac{1}{\sin\alpha\cos\alpha}\) en \(\tan\alpha+\cot\alpha\):
\[
\frac{1}{\sin\alpha\cos\alpha}
= \frac{\sin^2\alpha+\cos^2\alpha}{\sin\alpha\cos\alpha}
= \frac{\sin\alpha}{\cos\alpha} + \frac{\cos\alpha}{\sin\alpha}
= \tan\alpha + \cot\alpha
\]
Por lo tanto \(L = \tan\alpha+\cot\alpha\), que es el lado derecho.

\newpage
% =====================================================
% c)
% =====================================================
\section*{c)}

\begin{identidad}
\[
\frac{\sin\alpha}{1+\cos\alpha}
+
\cot\alpha
=
\csc\alpha
\]
\end{identidad}

\subsection*{Paso 1}
Partimos del lado izquierdo:
\[
L = \frac{\sin\alpha}{1+\cos\alpha} + \cot\alpha
= \frac{\sin\alpha}{1+\cos\alpha} + \frac{\cos\alpha}{\sin\alpha}
\]

\subsection*{Paso 2}
Sumamos las fracciones:
\[
L = \frac{\sin^2\alpha + \cos\alpha(1+\cos\alpha)}{\sin\alpha(1+\cos\alpha)}
= \frac{\sin^2\alpha + \cos\alpha + \cos^2\alpha}{\sin\alpha(1+\cos\alpha)}
\]

\subsection*{Paso 3}
Usamos \(\sin^2\alpha+\cos^2\alpha = 1\):
\[
L = \frac{1 + \cos\alpha}{\sin\alpha(1+\cos\alpha)}
\]
Cancelamos \(1+\cos\alpha\):
\[
L = \frac{1}{\sin\alpha} = \csc\alpha
\]

\newpage
% =====================================================
% d)
% =====================================================
\section*{d)}

\begin{identidad}
\[
\frac{\sec\alpha+1}{\sec\alpha-1}
-
\frac{\sec\alpha-1}{\sec\alpha+1}
-
4\cot^2\alpha
=
\frac4{1+\sec\alpha}
\]
\end{identidad}

\subsection*{Paso 1}
Partimos del lado izquierdo:
\[
L = \frac{\sec\alpha+1}{\sec\alpha-1} - \frac{\sec\alpha-1}{\sec\alpha+1} - 4\cot^2\alpha
\]
Restamos las dos primeras fracciones:
\[
\frac{\sec\alpha+1}{\sec\alpha-1} - \frac{\sec\alpha-1}{\sec\alpha+1}
= \frac{(\sec\alpha+1)^2 - (\sec\alpha-1)^2}{(\sec\alpha-1)(\sec\alpha+1)}
\]
Usando \((a+b)^2-(a-b)^2 = 4ab\) con \(a=\sec\alpha, b=1\):
\[
= \frac{4\sec\alpha}{\sec^2\alpha-1}
\]

\subsection*{Paso 2}
Sabemos \(\sec^2\alpha-1 = \tan^2\alpha\). Luego:
\[
\frac{4\sec\alpha}{\tan^2\alpha} = 4\sec\alpha \cdot \frac{\cos^2\alpha}{\sin^2\alpha}
= \frac{4}{\cos\alpha} \cdot \frac{\cos^2\alpha}{\sin^2\alpha}
= \frac{4\cos\alpha}{\sin^2\alpha}
\]

\subsection*{Paso 3}
Ahora restamos \(4\cot^2\alpha = \frac{4\cos^2\alpha}{\sin^2\alpha}\):
\[
L = \frac{4\cos\alpha}{\sin^2\alpha} - \frac{4\cos^2\alpha}{\sin^2\alpha}
= \frac{4\cos\alpha(1-\cos\alpha)}{\sin^2\alpha}
\]

\subsection*{Paso 4}
Usamos \(\sin^2\alpha = (1-\cos\alpha)(1+\cos\alpha)\):
\[
L = \frac{4\cos\alpha(1-\cos\alpha)}{(1-\cos\alpha)(1+\cos\alpha)}
= \frac{4\cos\alpha}{1+\cos\alpha}
\]

\subsection*{Paso 5}
Expresamos \(\frac{4\cos\alpha}{1+\cos\alpha}\) en función de secante:
\[
\frac{4\cos\alpha}{1+\cos\alpha}
= \frac{4}{ \frac{1}{\cos\alpha} + 1}
= \frac{4}{\sec\alpha+1}
\]
Hemos llegado al lado derecho.

\newpage
% =====================================================
% e)
% =====================================================
\section*{e)}

\begin{identidad}
\[
\frac{1+\tan^2\alpha}
{1+\cot^2\alpha}
=
\left(
\frac{1-\tan\alpha}
{1-\cot\alpha}
\right)^2
\]
\end{identidad}

\subsection*{Paso 1}
Partimos del lado izquierdo:
\[
L = \frac{1+\tan^2\alpha}{1+\cot^2\alpha}
= \frac{\sec^2\alpha}{\csc^2\alpha}
= \frac{1/\cos^2\alpha}{1/\sin^2\alpha}
= \frac{\sin^2\alpha}{\cos^2\alpha}
= \tan^2\alpha
\]

\subsection*{Paso 2}
Transformamos \(\tan^2\alpha\) en la forma del lado derecho. Sea \(t = \tan\alpha\). Entonces:
\[
\tan^2\alpha = t^2 = (-t)^2 = \left( (1-t)\cdot\frac{t}{t-1} \right)^2
\]
Notamos que \(\frac{t}{t-1} = \frac{1}{1 - 1/t} = \frac{1}{1-\cot\alpha}\), y \(1-t = -(t-1)\). Por tanto:
\[
-t = \frac{1-t}{1 - \frac1t} = \frac{1-\tan\alpha}{1-\cot\alpha}
\]
Luego:
\[
\tan^2\alpha = \left( \frac{1-\tan\alpha}{1-\cot\alpha} \right)^2
\]
Así \(L\) coincide exactamente con el lado derecho.

\newpage
% =====================================================
% f)
% =====================================================
\section*{f)}

\begin{identidad}
\[
\frac{
1-\sin\alpha\cos\alpha
}{
\cos\alpha(\sec\alpha-\csc\alpha)
}
\cdot
\frac{
\sin^2\alpha-\cos^2\alpha
}{
\sin^3\alpha+\cos^3\alpha
}
=
\sin\alpha
\]
\end{identidad}

\subsection*{Paso 1}
Partimos del producto del lado izquierdo. Simplificamos \(\cos\alpha(\sec\alpha-\csc\alpha)\):
\[
\sec\alpha-\csc\alpha = \frac{1}{\cos\alpha}-\frac{1}{\sin\alpha}
= \frac{\sin\alpha-\cos\alpha}{\sin\alpha\cos\alpha}
\]
Multiplicando por \(\cos\alpha\):
\[
\cos\alpha(\sec\alpha-\csc\alpha) = \frac{\sin\alpha-\cos\alpha}{\sin\alpha}
\]
Luego la primera fracción se convierte en:
\[
\frac{1-\sin\alpha\cos\alpha}{\frac{\sin\alpha-\cos\alpha}{\sin\alpha}}
= \frac{\sin\alpha(1-\sin\alpha\cos\alpha)}{\sin\alpha-\cos\alpha}
\]

\subsection*{Paso 2}
La segunda fracción:
\[
\sin^2\alpha-\cos^2\alpha = (\sin\alpha-\cos\alpha)(\sin\alpha+\cos\alpha)
\]
\[
\sin^3\alpha+\cos^3\alpha = (\sin\alpha+\cos\alpha)(\sin^2\alpha-\sin\alpha\cos\alpha+\cos^2\alpha)
= (\sin\alpha+\cos\alpha)(1-\sin\alpha\cos\alpha)
\]
Por tanto:
\[
\frac{\sin^2\alpha-\cos^2\alpha}{\sin^3\alpha+\cos^3\alpha}
= \frac{(\sin\alpha-\cos\alpha)(\sin\alpha+\cos\alpha)}{(\sin\alpha+\cos\alpha)(1-\sin\alpha\cos\alpha)}
= \frac{\sin\alpha-\cos\alpha}{1-\sin\alpha\cos\alpha}
\]

\subsection*{Paso 3}
Multiplicamos ambas fracciones:
\[
L = \frac{\sin\alpha(1-\sin\alpha\cos\alpha)}{\sin\alpha-\cos\alpha}
\cdot
\frac{\sin\alpha-\cos\alpha}{1-\sin\alpha\cos\alpha}
= \sin\alpha
\]
Lado derecho obtenido.

\newpage
% =====================================================
% g)
% =====================================================
\section*{g)}

\begin{identidad}
\[
2\sec^2\alpha-\sec^4\alpha-2\cot^2\alpha+\cot^4\alpha
=
\cot^4\alpha-\tan^4\alpha-2\cot^2\alpha+1
\]
\end{identidad}

\subsection*{Paso 1}
Partimos del lado izquierdo:
\[
L = 2\sec^2\alpha - \sec^4\alpha - 2\cot^2\alpha + \cot^4\alpha
\]
Escribimos \(\sec^2\alpha = 1+\tan^2\alpha\). Entonces:
\[
2\sec^2\alpha = 2(1+\tan^2\alpha) = 2+2\tan^2\alpha
\]
\[
\sec^4\alpha = (1+\tan^2\alpha)^2 = 1 + 2\tan^2\alpha + \tan^4\alpha
\]
Sustituimos:
\[
L = (2+2\tan^2\alpha) - (1+2\tan^2\alpha+\tan^4\alpha) - 2\cot^2\alpha + \cot^4\alpha
\]

\subsection*{Paso 2}
Simplificamos:
\[
L = 2+2\tan^2\alpha -1 -2\tan^2\alpha - \tan^4\alpha - 2\cot^2\alpha + \cot^4\alpha
\]
Los términos \(2\tan^2\alpha\) y \(-2\tan^2\alpha\) se cancelan:
\[
L = 1 - \tan^4\alpha - 2\cot^2\alpha + \cot^4\alpha
\]

\subsection*{Paso 3}
Reordenamos:
\[
L = \cot^4\alpha - \tan^4\alpha - 2\cot^2\alpha + 1
\]
Que es exactamente el lado derecho.

\newpage
% =====================================================
% h)
% =====================================================
\section*{h)}

\begin{identidad}
\[
\frac1{1+\sin^2\alpha}
+
\frac1{1+\cos^2\alpha}
=
\frac3{2+\sin^2\alpha\cos^2\alpha}
\]
\end{identidad}

\subsection*{Paso 1}
Partimos del lado izquierdo:
\[
L = \frac1{1+\sin^2\alpha} + \frac1{1+\cos^2\alpha}
= \frac{(1+\cos^2\alpha)+(1+\sin^2\alpha)}{(1+\sin^2\alpha)(1+\cos^2\alpha)}
\]

\subsection*{Paso 2}
Simplificamos el numerador:
\[
1+1 + \sin^2\alpha+\cos^2\alpha = 2 + 1 = 3
\]
Denominador:
\[
(1+\sin^2\alpha)(1+\cos^2\alpha) = 1 + \cos^2\alpha + \sin^2\alpha + \sin^2\alpha\cos^2\alpha
= 1 + 1 + \sin^2\alpha\cos^2\alpha = 2 + \sin^2\alpha\cos^2\alpha
\]
Por tanto:
\[
L = \frac{3}{2+\sin^2\alpha\cos^2\alpha}
\]

\newpage
% =====================================================
% i)
% =====================================================
\section*{i)}

\begin{identidad}
\[
\frac{\sin^6\alpha-\cos^6\alpha}
{\sin^2\alpha-\cos^2\alpha}
+
\sin^2\alpha\cos^2\alpha
=
1
\]
\end{identidad}

\subsection*{Paso 1}
Partimos del lado izquierdo:
\[
L = \frac{\sin^6\alpha-\cos^6\alpha}{\sin^2\alpha-\cos^2\alpha} + \sin^2\alpha\cos^2\alpha
\]
Factorizamos la diferencia de cubos:
\[
\sin^6\alpha-\cos^6\alpha = (\sin^2\alpha)^3 - (\cos^2\alpha)^3
= (\sin^2\alpha-\cos^2\alpha)(\sin^4\alpha + \sin^2\alpha\cos^2\alpha + \cos^4\alpha)
\]

\subsection*{Paso 2}
Cancelamos \(\sin^2\alpha-\cos^2\alpha\):
\[
L = \sin^4\alpha + \sin^2\alpha\cos^2\alpha + \cos^4\alpha + \sin^2\alpha\cos^2\alpha
= \sin^4\alpha + 2\sin^2\alpha\cos^2\alpha + \cos^4\alpha
\]

\subsection*{Paso 3}
Reconocemos un cuadrado perfecto:
\[
L = (\sin^2\alpha + \cos^2\alpha)^2 = 1^2 = 1
\]

\newpage
% =====================================================
% j)
% =====================================================
\section*{j)}

\begin{identidad}
\[
\frac{\tan^3\alpha}{1+\tan^2\alpha}
+
\frac{\cot^3\alpha}{1+\cot^2\alpha}
=
\sec\alpha\csc\alpha-2\sin\alpha\cos\alpha
\]
\end{identidad}

\subsection*{Paso 1}
Partimos del lado izquierdo:
\[
L = \frac{\tan^3\alpha}{1+\tan^2\alpha} + \frac{\cot^3\alpha}{1+\cot^2\alpha}
\]
Usamos \(1+\tan^2\alpha = \sec^2\alpha\) y \(1+\cot^2\alpha = \csc^2\alpha\):
\[
L = \frac{\tan^3\alpha}{\sec^2\alpha} + \frac{\cot^3\alpha}{\csc^2\alpha}
\]

\subsection*{Paso 2}
Escribimos en seno y coseno:
\[
\frac{\tan^3\alpha}{\sec^2\alpha} = \frac{\frac{\sin^3\alpha}{\cos^3\alpha}}{\frac1{\cos^2\alpha}} = \frac{\sin^3\alpha}{\cos\alpha}
\]
\[
\frac{\cot^3\alpha}{\csc^2\alpha} = \frac{\frac{\cos^3\alpha}{\sin^3\alpha}}{\frac1{\sin^2\alpha}} = \frac{\cos^3\alpha}{\sin\alpha}
\]
Por tanto:
\[
L = \frac{\sin^3\alpha}{\cos\alpha} + \frac{\cos^3\alpha}{\sin\alpha}
= \frac{\sin^4\alpha + \cos^4\alpha}{\sin\alpha\cos\alpha}
\]

\subsection*{Paso 3}
Usamos \(\sin^4\alpha+\cos^4\alpha = (\sin^2\alpha+\cos^2\alpha)^2 - 2\sin^2\alpha\cos^2\alpha = 1 - 2\sin^2\alpha\cos^2\alpha\):
\[
L = \frac{1 - 2\sin^2\alpha\cos^2\alpha}{\sin\alpha\cos\alpha}
= \frac{1}{\sin\alpha\cos\alpha} - 2\sin\alpha\cos\alpha
\]

\subsection*{Paso 4}
Finalmente:
\[
\frac{1}{\sin\alpha\cos\alpha} = \csc\alpha\sec\alpha
\]
Luego:
\[
L = \csc\alpha\sec\alpha - 2\sin\alpha\cos\alpha
\]

\newpage
% =====================================================
% k)
% =====================================================
\section*{k)}

\begin{identidad}
\[
\sec\theta-\cos\theta=\tan\theta\sin\theta
\]
\end{identidad}

\subsection*{Paso 1}
Partimos del lado izquierdo:
\[
L = \sec\theta - \cos\theta = \frac1{\cos\theta} - \cos\theta
\]

\subsection*{Paso 2}
Común denominador:
\[
L = \frac{1-\cos^2\theta}{\cos\theta} = \frac{\sin^2\theta}{\cos\theta}
\]

\subsection*{Paso 3}
Separamos:
\[
L = \sin\theta \cdot \frac{\sin\theta}{\cos\theta} = \sin\theta\tan\theta
\]

\newpage
% =====================================================
% l)
% =====================================================
\section*{l)}

\begin{identidad}
\[
\csc\theta-\sin\theta=\cot\theta\cos\theta
\]
\end{identidad}

\subsection*{Paso 1}
Partimos del lado izquierdo:
\[
L = \csc\theta - \sin\theta = \frac1{\sin\theta} - \sin\theta
\]

\subsection*{Paso 2}
Común denominador:
\[
L = \frac{1-\sin^2\theta}{\sin\theta} = \frac{\cos^2\theta}{\sin\theta}
\]

\subsection*{Paso 3}
\[
L = \cos\theta \cdot \frac{\cos\theta}{\sin\theta} = \cos\theta\cot\theta
\]

\newpage
% =====================================================
% m)
% =====================================================
\section*{m)}

\begin{identidad}
\[
\frac{\sin\theta+\cos\theta}{\sin\theta}
=
1+\cot\theta
\]
\end{identidad}

\subsection*{Paso 1}
Partimos del lado izquierdo:
\[
L = \frac{\sin\theta+\cos\theta}{\sin\theta}
= \frac{\sin\theta}{\sin\theta} + \frac{\cos\theta}{\sin\theta}
= 1 + \cot\theta
\]

\newpage
% =====================================================
% n)
% =====================================================
\section*{n)}

\begin{identidad}
\[
\frac{\sin\theta+\cos\theta}{\cos\theta}
=
1+\tan\theta
\]
\end{identidad}

\subsection*{Paso 1}
Partimos del lado izquierdo:
\[
L = \frac{\sin\theta+\cos\theta}{\cos\theta}
= \frac{\sin\theta}{\cos\theta} + \frac{\cos\theta}{\cos\theta}
= \tan\theta + 1
\]

\newpage
% =====================================================
% ñ)
% =====================================================
\section*{ñ)}

\begin{identidad}
\[
(\cot\theta+\csc\theta)(\tan\theta-\sin\theta)
=
\sec\theta-\cos\theta
\]
\end{identidad}

\subsection*{Paso 1}
Partimos del lado izquierdo:
\[
L = (\cot\theta+\csc\theta)(\tan\theta-\sin\theta)
\]
Escribimos cada factor en seno y coseno:
\[
\cot\theta+\csc\theta = \frac{\cos\theta}{\sin\theta} + \frac1{\sin\theta} = \frac{\cos\theta+1}{\sin\theta}
\]
\[
\tan\theta-\sin\theta = \frac{\sin\theta}{\cos\theta} - \sin\theta = \sin\theta\left(\frac1{\cos\theta}-1\right)
\]

\subsection*{Paso 2}
Multiplicamos:
\[
L = \frac{\cos\theta+1}{\sin\theta} \cdot \sin\theta\left(\frac1{\cos\theta}-1\right)
= (\cos\theta+1)\left(\frac1{\cos\theta}-1\right)
\]

\subsection*{Paso 3}
Simplificamos:
\[
L = (\cos\theta+1)\cdot\frac{1-\cos\theta}{\cos\theta}
= \frac{(1+\cos\theta)(1-\cos\theta)}{\cos\theta}
= \frac{1-\cos^2\theta}{\cos\theta}
= \frac{\sin^2\theta}{\cos\theta}
\]

\subsection*{Paso 4}
\[
L = \frac1{\cos\theta} - \cos\theta = \sec\theta - \cos\theta
\]

\newpage
% =====================================================
% o)
% =====================================================
\section*{o)}

\begin{identidad}
\[
\cot\theta+\tan\theta
=
\csc\theta\sec\theta
\]
\end{identidad}

\subsection*{Paso 1}
Partimos del lado izquierdo:
\[
L = \cot\theta+\tan\theta = \frac{\cos\theta}{\sin\theta} + \frac{\sin\theta}{\cos\theta}
\]

\subsection*{Paso 2}
Sumamos:
\[
L = \frac{\cos^2\theta+\sin^2\theta}{\sin\theta\cos\theta}
= \frac{1}{\sin\theta\cos\theta}
\]

\subsection*{Paso 3}
\[
L = \frac1{\sin\theta} \cdot \frac1{\cos\theta} = \csc\theta\sec\theta
\]

\newpage
% =====================================================
% p)
% =====================================================
\section*{p)}

\begin{identidad}
\[
\sec^2(3\theta)\csc^2(3\theta)
=
\sec^2(3\theta)+\csc^2(3\theta)
\]
\end{identidad}

\subsection*{Paso 1}
Partimos del lado izquierdo:
\[
L = \sec^2(3\theta)\csc^2(3\theta) = \frac{1}{\cos^2(3\theta)} \cdot \frac{1}{\sin^2(3\theta)}
= \frac{1}{\sin^2(3\theta)\cos^2(3\theta)}
\]

\subsection*{Paso 2}
Transformamos esta expresión en la suma:
\[
\frac{1}{\sin^2(3\theta)\cos^2(3\theta)}
= \frac{\sin^2(3\theta)+\cos^2(3\theta)}{\sin^2(3\theta)\cos^2(3\theta)}
= \frac{1}{\cos^2(3\theta)} + \frac{1}{\sin^2(3\theta)}
= \sec^2(3\theta) + \csc^2(3\theta)
\]
Hemos obtenido el lado derecho.

\newpage
% =====================================================
% q)
% =====================================================
\section*{q)}

\begin{identidad}
\[
\frac{1+\cos^2(3\theta)}
{\sin^2(3\theta)}
=
2\csc^2(3\theta)-1
\]
\end{identidad}

\subsection*{Paso 1}
Partimos del lado izquierdo:
\[
L = \frac{1+\cos^2(3\theta)}{\sin^2(3\theta)}
\]
Usamos \(\cos^2(3\theta) = 1 - \sin^2(3\theta)\):
\[
L = \frac{1 + 1 - \sin^2(3\theta)}{\sin^2(3\theta)}
= \frac{2 - \sin^2(3\theta)}{\sin^2(3\theta)}
\]

\subsection*{Paso 2}
Separamos:
\[
L = \frac{2}{\sin^2(3\theta)} - 1 = 2\csc^2(3\theta) - 1
\]

\newpage
% =====================================================
% r)
% =====================================================
\section*{r)}

\begin{identidad}
\[
\log(\csc\theta)
=
-\log(\sin\theta)
\]
\end{identidad}

\subsection*{Paso 1}
Partimos del lado izquierdo:
\[
L = \log(\csc\theta) = \log\left(\frac{1}{\sin\theta}\right)
\]

\subsection*{Paso 2}
Usamos la propiedad de logaritmos:
\[
\log\left(\frac{1}{\sin\theta}\right) = \log(1) - \log(\sin\theta) = 0 - \log(\sin\theta) = -\log(\sin\theta)
\]

\end{document}